%! Intercepts, Zeros, and Extrema — Unit 2, Lesson 2.4 — Guided Notes
\documentclass[10pt]{article}
\usepackage{atda-article}
\usepackage{atda-boxes}
\newcommand{\TallMath}[1]{$\displaystyle #1\rule[-1.4em]{0pt}{3.2em}$}
\setlength{\workrowsep}{6pt}

\begin{document}

\pageheader{Unit 2, Lesson 2.4}{Guided Notes}
% NAMESTRIP: no \namedateperiod here. The student writes their name once, on the
% cover this component is stapled behind. See references/conventions.md.

\vspace{0.15in}

\begin{objectivebox}
\small
By the end of this lesson, I will be able to\dots
\begin{enumerate}[leftmargin=1.5em, itemsep=2pt, topsep=3pt]
  \item find the \textbf{$y$-intercept} and the \textbf{$x$-intercepts} of a graph and give the
        \textbf{zeros} of the function --- saying each time whether my answer is an input, an output,
        or a point;
  \item name the intervals where a graph is \textbf{increasing}, \textbf{decreasing}, or
        \textbf{constant}, and write them in interval notation using the \emph{input} axis;
  \item identify the \textbf{absolute maximum} and \textbf{absolute minimum} over the domain,
        reporting a \emph{value} and a \emph{location} every time;
  \item explain why a function may have no zeros, no absolute maximum, or two places where the same
        extreme value happens --- and interpret every characteristic in context.
\end{enumerate}
\end{objectivebox}

\vspace{0.10in}

\boxguard
\begin{vocabbox}
\small
Fill in each term as we build it together. Every one of them is read straight off a graph.
\termblanklong{Zero of a function}
\termblanklong{$x$-intercept and $y$-intercept}
\termblanklong{Increasing, decreasing, and constant}
\termblanklong{Turning point}
\termblanklong{Absolute maximum and absolute minimum}
\end{vocabbox}

\vspace{0.10in}

\boxguard[12]
\begin{hookbox}
\small
\textbf{Two students, one question, and both answers are in dollars.} Last night you worked with the
yearbook's profit curve: profit is \$0 at prices of \$0 and \$60, and the profit climbs to a high of
\$1800 at a price of \$30. The question was: \emph{over which prices is the profit going up?}

\smallskip
Student A wrote \quad \textbf{``from $0$ to $1800$''} \qquad Student B wrote \quad
\textbf{``from $0$ to $30$''}

\smallskip
Neither of them made an arithmetic mistake, and every number either one wrote is really on that
graph. \textbf{Here is the question: which one answered the question that was asked, and how would
you convince the other one?} Notice that you cannot settle it by checking the units --- all four
numbers are dollars.

\writelines{2}
\end{hookbox}

\vspace{0.10in}

\boxguard[30]
\begin{notesbox}{1. Intercepts and Zeros}
\small
The first thing you look for on any graph is where it meets the axes. There are only two such
questions, and each one is really a question about a \emph{coordinate that is zero}.

\vspace{4pt}
\textbf{The $y$-intercept} is the point where the graph crosses the \blank{2.4cm} axis. Getting
there means the input is $x=$ \blank{0.9cm}. A function has \emph{at most one} $y$-intercept --- one
input, one output.

\vspace{2pt}
\textbf{An $x$-intercept} is a point where the graph crosses the \blank{2.4cm} axis. Being there
means the output is $y=$ \blank{0.9cm}. A graph may have one, several, or \emph{none at all}.

\vspace{2pt}
\textbf{A \blank{2.2cm}} is the \emph{input} at an $x$-intercept. So the $x$-intercept is a
\textbf{point} and the zero is the \textbf{number} in its first slot: if $(5,0)$ is on the graph,
then $5$ is a zero. These three sentences say the same thing:
\begin{itemize}[leftmargin=1.3em, itemsep=1pt, topsep=3pt]
  \item $k$ is a \textbf{zero} of $f$; \quad $f(k)=0$; \quad the point $(k,0)$ is on the graph of $f$.
\end{itemize}

\vspace{4pt}
\begin{center}
\begin{tabular}{@{}c@{\hspace{0.4cm}}c@{\hspace{0.4cm}}c@{}}
\begin{tikzpicture}
\begin{axis}[
    width=5.0cm, height=4.2cm, title={\scriptsize (a) $y=x^2-4$},
    xmin=-3.4, xmax=3.4, ymin=-5.4, ymax=5.4,
    xtick={-3,-2,-1,0,1,2,3}, ytick={-4,-2,0,2,4}, minor y tick num=1,
    grid=both, grid style={linegray!45}, minor grid style={linegray!30},
    axis lines=left, tick label style={font=\tiny}]
  \addplot[royal, very thick, domain=-3:3, samples=90]{x^2-4};
\end{axis}
\end{tikzpicture}
&
\begin{tikzpicture}
\begin{axis}[
    width=5.0cm, height=4.2cm, title={\scriptsize (b) $y=2^x$},
    xmin=-3.4, xmax=3.4, ymin=-1.4, ymax=8.4,
    xtick={-3,-2,-1,0,1,2,3}, ytick={0,2,4,6,8}, minor y tick num=1,
    grid=both, grid style={linegray!45}, minor grid style={linegray!30},
    axis lines=left, tick label style={font=\tiny}]
  \addplot[royal, very thick, domain=-3.3:3, samples=90]{2^x};
\end{axis}
\end{tikzpicture}
&
\begin{tikzpicture}
\begin{axis}[
    width=5.0cm, height=4.2cm, title={\scriptsize (c) $y=3-x$},
    xmin=-1.4, xmax=5.4, ymin=-2.4, ymax=4.4,
    xtick={-1,0,1,2,3,4,5}, ytick={-2,0,2,4}, minor y tick num=1,
    grid=both, grid style={linegray!45}, minor grid style={linegray!30},
    axis lines=left, tick label style={font=\tiny}]
  \addplot[royal, very thick, domain=-1.3:5.3, samples=2]{3-x};
\end{axis}
\end{tikzpicture}
\end{tabular}
\end{center}

\vspace{-2pt}
\renewcommand{\arraystretch}{1.5}
% Headers break explicitly: at \small the one-line form "y-intercept (a point)"
% overran its column and wrapped mid-parenthesis. Mirrored in notes_key.
\begin{tabularx}{\linewidth}{@{}p{1.5cm} p{3.4cm} p{4.4cm} X@{}}
\rowcolor{mist}
\textbf{Graph} & \textbf{$y$-intercept}\newline\textbf{(a point)} &
    \textbf{$x$-intercept(s)}\newline\textbf{(points)} &
    \textbf{Zero(s)}\newline\textbf{(numbers)} \\
\hline
(a) & \blank{2.4cm} & \blank{3.4cm} & \blank{2.8cm} \\
(b) & \blank{2.4cm} & \blank{3.4cm} & \blank{2.8cm} \\
(c) & \blank{2.4cm} & \blank{3.4cm} & \blank{2.8cm} \\
\end{tabularx}

\vspace{6pt}
\textbf{Two things worth writing down.} Graph (b) has \emph{no} zeros: the curve gets closer and
closer to the horizontal axis on the left, but it never arrives. \textbf{Not every graph crosses the
$x$-axis.} And in graph (c) the $y$-intercept and the zero are both the number $3$, which is a
coincidence and nothing more --- one of them is an \blank{1.9cm} and the other is an
\blank{1.9cm}. \textbf{Before you write a number down, say which axis it came from.}
\end{notesbox}

\vspace{0.10in}

% Pagination verdict, measured: notes runs 5 pages with each of pp. 2--5 about 60%
% full, and that is the natural break, not slack to reclaim. Each of the four boxes
% is taller than half a page (box 1 ~5.5in, box 2 ~6.5in, box 3 ~6in, practice
% ~5.5in against ~9in of column), so no two can share a page. Lowering this guard to
% let box 2 start on p2 was measured and rejected: p2 has ~2.5in free and box 2's
% first unbreakable chunk (intro line + the 5.4cm axis) is ~2.9in, so it would strand
% the intro line alone. Restored per 1.5's test-then-restore rule.
\boxguard[30]
\begin{notesbox}{2. Rising, Falling, and Holding Steady}
\small
A graph is read \textbf{left to right} --- always, in the direction the \blank{2.2cm} increases. As
you sweep across, the curve is doing one of exactly three things.

\begin{center}
\begin{tikzpicture}
\begin{axis}[
    width=11.0cm, height=5.4cm,
    xlabel={$t$ (hours after midnight)}, ylabel={$T$ (degrees C)},
    xmin=0, xmax=12, ymin=-6.4, ymax=6.4,
    xtick={0,2,4,6,8,10,12}, ytick={-6,-4,-2,0,2,4,6},
    minor x tick num=1, minor y tick num=1,
    grid=both, grid style={linegray!45}, minor grid style={linegray!30},
    axis lines=left,
    tick label style={font=\tiny}, label style={font=\scriptsize}]
  \addplot[royal, very thick] coordinates {(0,2) (2,2) (5,-4) (9,4) (12,1)};
\end{axis}
\end{tikzpicture}
\end{center}

\vspace{-4pt}
{\scriptsize A cold front moved through overnight and clouds rolled in mid-morning. $T(t)$ is the
temperature in degrees Celsius, $t$ hours after midnight, for the twelve hours to noon.}

\vspace{4pt}
\textbf{Fill the table from the graph. Every interval is a set of \emph{times}.}

\vspace{2pt}
\renewcommand{\arraystretch}{1.5}
\begin{tabularx}{\linewidth}{@{}p{2.4cm} p{3.2cm} X@{}}
\rowcolor{mist}
\textbf{Behavior} & \textbf{On these inputs} & \textbf{What it means here} \\
\hline
Constant & \blank{2.4cm} & \blank{6.0cm} \\
Decreasing & \blank{2.4cm} & \blank{6.0cm} \\
Increasing & \blank{2.4cm} & \blank{6.0cm} \\
Decreasing & \blank{2.4cm} & \blank{6.0cm} \\
\end{tabularx}

\vspace{6pt}
\textbf{Three rules that decide every one of these.}
\begin{enumerate}[leftmargin=1.5em, itemsep=3pt, topsep=3pt]
  \item \textbf{Report the interval on the \blank{2.6cm} axis.} Increasing, decreasing and constant
        are answers about \emph{which inputs}, never about how high the graph got. This is the whole
        hook: student \blank{0.9cm} answered the question that was asked, because prices are the
        input. Student A described the outputs the graph swept through --- true, and not the
        question.
  \item \textbf{Use parentheses at the ends.} At $t=5$ the graph is neither falling nor rising --- it
        is turning. A point where the graph changes from decreasing to increasing (or the other way)
        is a \textbf{turning point}, and we leave it out of both intervals.
  \item \textbf{An exponential never changes its mind.} $y=2^x$ rises across its entire domain and
        $y=\left(\tfrac12\right)^x$ falls across its entire domain. An exponential is
        \blank{2.4cm} increasing or \blank{2.4cm} decreasing --- one interval, and it is all of
        $(-\infty,\infty)$.
\end{enumerate}

\vspace{4pt}
\textbf{And read the zeros back into the story.} This graph has zeros at $t=3$ and $t=7$. Between
them the curve is below the axis, so the temperature was below freezing from $3$~a.m.\ to
$7$~a.m. \textbf{Two zeros gave you an interval of time} --- that is what zeros are for.
\end{notesbox}

\vspace{0.10in}

\boxguard[30]
\begin{notesbox}{3. The Highest and the Lowest --- and Why the Domain Decides}
\small
The \textbf{absolute maximum} of a function is its highest output over the \emph{entire domain};
the \textbf{absolute minimum} is its lowest. Both answers come in two pieces.

\vspace{4pt}
\textbf{Report a \blank{1.8cm} and a \blank{1.8cm}.} ``The maximum is $4$ degrees'' is half an
answer; ``$4$ degrees, at $t=9$'' is the whole one. The value is read on the vertical axis and the
location on the horizontal axis --- \emph{two different axes in one sentence}, which is exactly why
this is where people slip.

\vspace{2pt}
\textbf{``Absolute'' means over the \blank{3.0cm}.} A high point in the middle of a graph is not
automatically the absolute maximum; you have to look at the whole picture, endpoints included.

\vspace{4pt}
\begin{center}
\begin{tabular}{@{}c@{\hspace{0.8cm}}c@{}}
\begin{tikzpicture}
\begin{axis}[
    width=6.2cm, height=4.4cm, title={\scriptsize (a) $y=x^2$, domain all real numbers},
    xmin=-3.4, xmax=3.4, ymin=-1.0, ymax=10.4,
    xtick={-3,-2,-1,0,1,2,3}, ytick={0,2,4,6,8,10}, minor y tick num=1,
    grid=both, grid style={linegray!45}, minor grid style={linegray!30},
    axis lines=left, tick label style={font=\tiny}]
  \addplot[royal, very thick, domain=-3.15:3.15, samples=90]{x^2};
\end{axis}
\end{tikzpicture}
&
\begin{tikzpicture}
\begin{axis}[
    width=6.2cm, height=4.4cm, title={\scriptsize (b) $y=x^2$, domain $-1 \le x \le 3$},
    xmin=-3.4, xmax=3.4, ymin=-1.0, ymax=10.4,
    xtick={-3,-2,-1,0,1,2,3}, ytick={0,2,4,6,8,10}, minor y tick num=1,
    grid=both, grid style={linegray!45}, minor grid style={linegray!30},
    axis lines=left, tick label style={font=\tiny}]
  \addplot[royal, very thick, domain=-1:3, samples=90]{x^2};
  \addplot[royal, only marks, mark=*, mark size=1.7pt] coordinates {(-1,1) (3,9)};
\end{axis}
\end{tikzpicture}
\end{tabular}
\end{center}

\vspace{-2pt}
{\scriptsize The same rule, $y=x^2$, on two different domains. Nothing about the formula changed.}

\vspace{4pt}
\renewcommand{\arraystretch}{1.5}
\begin{tabularx}{\linewidth}{@{}p{4.2cm} X X@{}}
\rowcolor{mist}
 & \textbf{Absolute minimum} & \textbf{Absolute maximum} \\
\hline
(a) domain all real numbers & \blank{3.6cm} & \blank{3.6cm} \\
(b) domain $-1 \le x \le 3$ & \blank{3.6cm} & \blank{3.6cm} \\
\end{tabularx}

\vspace{6pt}
\textbf{So the domain is part of the question, not decoration.} Cutting the domain down to
$-1 \le x \le 3$ did not move a single point of the curve, and it handed the function an absolute
maximum it did not have before --- at an \textbf{endpoint}, not at the turning point.

\vspace{4pt}
\textbf{Four cases to keep in your notes.}
\begin{itemize}[leftmargin=1.3em, itemsep=2pt, topsep=3pt]
  \item A parabola opening \textbf{upward} has an absolute minimum at its turning point and
        \emph{no} absolute maximum (unless the domain is cut).
  \item A parabola opening \textbf{downward} has an absolute maximum at its turning point and no
        absolute minimum.
  \item An exponential on its full domain has \emph{neither}: $y=2^x$ climbs forever, and going left
        it gets closer to $0$ without ever reaching it, so there is no lowest output.
  \item On a \textbf{restricted domain}, always check the \textbf{endpoints} --- an extreme is often
        sitting at one of them.
\end{itemize}
\end{notesbox}

\vspace{0.10in}

\boxguard[30]
\begin{practicebox}
\small
The school greenhouse collects rainwater in a cistern. The graph shows $W(t)$, the depth of the
water in feet, $t$ days after Monday morning, for one week.

\begin{center}
\begin{tikzpicture}
\begin{axis}[
    width=10.6cm, height=4.8cm,
    xlabel={$t$ (days after Monday morning)}, ylabel={$W$ (depth, feet)},
    xmin=0, xmax=7, ymin=0, ymax=8,
    xtick={0,1,2,3,4,5,6,7}, ytick={0,1,2,3,4,5,6,7,8},
    grid=both, grid style={linegray!45}, axis lines=left,
    tick label style={font=\tiny}, label style={font=\scriptsize}]
  \addplot[royal, very thick] coordinates {(0,6) (3,0) (5,4) (6,4) (7,7)};
\end{axis}
\end{tikzpicture}
\end{center}

\begin{enumerate}[leftmargin=1.6em, itemsep=6pt, topsep=2pt]
  \item Give the $y$-intercept as a point, and say what it means about the cistern.

\writelines{1}

  \item Give the zero of $W$, and say what it means about the cistern. Is it an input or an output?

\writelines{2}

  \item Name every interval on which $W$ is increasing, decreasing, or constant. Write each one in
        interval notation, and remember which axis they live on.

\writelines{3}

  \item Give the absolute maximum and the absolute minimum, each as a \emph{value and a location}.
        Then settle this claim: \emph{``The water was at its highest at the start of the week.''}

\writelines{3}
\end{enumerate}
\end{practicebox}

\end{document}
